\documentclass[11pt, reqno]{article}

\usepackage{amsmath, amsthm, amssymb}
\usepackage{enumitem}
\usepackage{tcolorbox}
\usepackage{hyperref}
\usepackage{tikz}
\usepackage{tikz-cd}
\usepackage{pgfplots}
\pgfplotsset{compat=1.18}
\usetikzlibrary{arrows.meta}
\usepackage{mathrsfs}
\usepackage{fancyhdr}
\usepackage[bottom=0.75in, top=1in, left=0.5in, right=0.5in]{geometry}
\usepackage{array}   % for \newcolumntype macro
\newcolumntype{L}{>{$}l<{$}}

\theoremstyle{plain}
\newtheorem*{theorem}{Theorem}
\newtheorem*{proposition}{Proposition}
\newtheorem{exercise}{Exercise}
\newtheorem*{lemma}{Lemma}
\newtheorem*{corollary}{Corollary}

\theoremstyle{definition}
\newtheorem*{definition}{Definition}
\newtheorem*{example}{Example}

\theoremstyle{remark}
\newtheorem*{remark}{Remark}

\renewcommand{\phi}{\varphi}
\renewcommand{\epsilon}{\varepsilon}
\renewcommand{\emptyset}{\varnothing}

\newcommand{\RR}{\mathbb{R}}
\newcommand{\ZZ}{\mathbb{Z}}
\newcommand{\NN}{\mathbb{N}}
\newcommand{\CC}{\mathbb{C}}
\newcommand{\QQ}{\mathbb{Q}}

\DeclareMathOperator{\ima}{\text{im}}

\begin{document}

\topmargin=-40pt
\rhead{Henry Woodburn}
\lhead{Math 622}
\renewcommand{\headrulewidth}{1pt}
\renewcommand{\headsep}{20pt}
\thispagestyle{fancy}

{\Huge \bfseries \noindent Homework 11}

\begin{enumerate}

    \item[1.] For $v \in T_g(G)$, we define $\omega_g(v) = d(L_{g^{-1}})(v)$. Then we can
    write 
    \begin{align*}
        (R_g)^*(\omega)_h (v) = \omega_{hg}(d(R_g)_h(v)) = 
        d(L_{g^{-1}h^{-1}}) \circ d(R_g)(v) = d(R_g) \circ d(L_{g^{-1}})\circ d(L_{h^{-1}})(v)\\
        = d(L_g \circ r_{g^{-1}})\omega(v),
    \end{align*}
    since it is clear that the differentials of left and right multiplication commute,
    since the underlying maps commute. 

    \item[2.] Define a distribution $D$ on $M \times N$ defined by the forms
    \[
        s_i = \pi^*(\alpha_i) - \rho^*(\omega_i).
    \]

    Then to show $D$ is involutive, we need to verify that $ds_i$ also vanishes on $D$. 
    We have
    \begin{align*}
        d(\pi^*\alpha_i - \rho^*\omega_i) = \pi^*(d\alpha_i) - \rho^*(d\omega_i) = 
        \pi^*(k_i \alpha_1 \wedge \wedge 2) - \rho^*(k_i \omega_1 \wedge \omega_2)\\
        = k_i(\pi^*(\alpha_1)\wedge \pi^*(\alpha_1) - \rho^*(\omega_1) \wedge \rho^*(\omega_2)).
    \end{align*}

    Then
    \[
        (\pi^*(\alpha_1) - \rho^*(\omega_1))\wedge(\pi^*(\alpha_2) + \rho^*(\omega_2))
        - (\pi^*(\alpha_2) - \rho^*(\omega_2))\wedge(\pi^*(\alpha_1) + \rho^*(\omega_1))
        = 2 \pi^*(\alpha_1)\wedge\pi^*(\alpha_2) - \rho^*(\omega_1)\wedge\rho^*(\omega_2),
    \]
    and from here it is clear that we can write
    \[
        ds_i = \sum_i \beta_i \wedge s_i
    \]
    by adding the constant $k_i/2$ to the previous calculation.

    Then $D$ is involutive. 

    Now we have for every $(m,n) \in M\times N$, there is a neighborhood $U \times V \ni (m,n)$
    and an integral submanifold $S$ defined on $U \times V$ passing through $(m,n)$,
    and a coordinate neighborhood such that all such integral manifolds are slices,
    with two coordinates constant. 

    Then since both projection maps are injections onto rank 2 manifolds, they are diffeomorphisms.
    Now define a diffeomorphism $F := \rho^{-1}\circ \pi$ on $U \subset N$. 

    We know that $\pi^*(\alpha_i) = \rho^*(\omega_i)$ on $S$, so that 
    \[
        \alpha_i = F^*(\omega_i).
    \]

    \item[3.] (a.) Suppose $u$ is a solution. Then $du(X_i) = X_i(u) = g_i = omega(X_i)$. Since
    the vector fields $X_i$ form a frame, it is enough to check this on each $X_i$. 
    Then $du = \omega$.

    Now suppose $du = \omega$. So $g_i = \omega(X_i) = du(X_i) = X_i(u)$.
    
    (b.) Suppose $u$ has a solution. Then $X_i(u) = g_i$. Then we have
    \[
        [X_1, X_2](u) = X_1(X_2(u)) - X_2(X_1(u)) = X_1(g_2) - X_2(g_1),
    \]
    and on the other hand we have
    \[
        [X_1, X_2](u) = \alpha_1 X_1(u) + \alpha_2 X_2(u) = \alpha_1 g_1 + \alpha_2 g_2.
    \]

    Now suppose the converse, that $X_1(g_2) - X_2(g_1) = \alpha_1 g_1 + \alpha_2 g_2$. We
    will show $\omega = du$ for some $u$, and thus $u$ is a solution by part (a.). To do 
    this we show $\omega$ is closed, which will imply it is exact since $\RR^2$ is contractible. 

    We use the coordinate free formula:
    \[
        d\omega(X_1, X_2) = X_1(\omega(X_2)) - X_2(\omega(X_1)) - \omega([X_1, X_2])
        = X_1(g_2) - X_2(g_1) - \omega(\alpha_1 X_1 + \alpha_2 X_2)
        = X_1 g_2 - X_2 g_1 - \alpha_1 g_1 - \alpha_2 g_2 = 0
    \]
    by assumption. 

    (ci.) Suppose $u$ is a solution. Choose $p = (x,y, u(x,y)) \in \Gamma_u$. Then since
    the vector fields form a frame, we can express the tangent space to $\Gamma_u$ at $p$
    as $\text{span}((1,0,\frac{\partial u}{\partial x_1}(x,y)), (0,1,\frac{\partial u}{\partial x_2}(x,y)))$,
    where $x_i = X_i(p)$. Then we check that 
    \[
        \Omega(1,0,\frac{\partial u}{\partial x_1}(x,y)) = \frac{\partial u}{\partial x_1}(x,y) - \omega(X_1)
        = \frac{\partial u}{\partial x_1}(x,y) - g_1(x,y).
    \]
    But the derivative of $u$ in the direction of $X_1$ at $(x,y)$ is just $g(x,y)$ since
    $u$ is a solution of (1). The same is true for the other basis vector. Then 
    is is clear the graph is integral. 

    Now suppose that the graph is integral. Then at any point, the derivatives in $X_1$ and $X_2$ 
    directions are annihilated by $\Omega$. Similar to the opposite direction, this 
    implies that $u$ is a solution to (1).

    \item[4.] Choose a coordinate neighborhood $U$ with coordinates $\{x_1, \dots, x_n\}$.
    Then the tautological symplectic 2-form is given by
    \[
        \omega = \sum_i dy_i \wedge dx^i.
    \]
    Then using the definition of the pullback, we have
    \[
        \iota^*(\omega) = \sum_i d(df(x)_i)\wedge d(x_i) = 0,
    \]

    since $df(x)_i = \frac{\partial f}{\partial x_i} dx_i$.

    \item[5.] Since $\theta$ is a contact form, there is a vector field $R$ such that 
    $\theta(R) = 1$ and $d\theta(R) = 0$. Then for any vector field $X$ we can write
    \[
        X = Y + fR
    \] 
    for a unique vector field $Y$ which is an element of $D$. This implies $d\theta(X) = d\theta(Y)$.
    Moreover, the condition that $X$ is an infinitesimal generator means 
    $\mathcal{L}_X\theta = g\theta$ for some function $g$. This is along the lines of the problem
    on a previous homework, since the distribution is generated a single one form. 

    We can choose a vector field $X_f$ such that $\theta(X_f) = f$, but it may not be
    an infinitesimal generator. 
    
    We have Cartan's Magic formula
    \[
        \mathcal{L}_X\theta = \iota_X(d\theta) + d(\iota_X\theta) = \iota_X(d\theta) + df.
    \]

    I could not continue.

    We can explicitly show that $X_f = x\frac{\partial}{\partial x} + y\frac{\partial}{\partial y} + f
    \frac{\partial}{\partial z}$.

    
\end{enumerate}

\end{document}